\section{Conclusion}
\label{sec:conclusion}

Result portals face a load that is announced in advance, read-only and simple,
yet they keep failing. We built ResultBench, which runs eleven portal designs
under the same two-CPU budget, replays a result-day surge with students who
refresh, and checks every answer against ground truth. On one machine, designs
that build pages on request failed by twice the capacity of a database-backed
portal, while pre-computed results served ten times that load at a small
fraction of the CPU. No design returned wrong results; stale results came only
from copies that were not updated after a correction. The weakness of
cache-aside on result day is cold misses, not stampedes, and refreshing
students can keep an unprotected portal overloaded.
A waiting room turned failures into long waits rather than serving more
students, and protected the backend only while admitted students behaved.

Several questions remain open. We have not measured a real result day; a field
study, for example by monitoring portals during a result season or by
obtaining load data from boards, would show how close our surge and student
model are to reality. Our page-building cost is a model; repeating the
comparison with real portal stacks, and with other values of that cost, would
test how far the gaps we measured carry over. Finally, all our runs used one
machine and no content delivery network; multi-machine and CDN deployments are
the next step.
